\documentclass[11pt]{article}

\usepackage[margin=1in]{geometry}
\usepackage{amsmath,amssymb,amsthm}
\usepackage{array}
\usepackage{booktabs}
\usepackage{longtable}
\usepackage{microtype}
\usepackage[hidelinks]{hyperref}

\newtheorem{theorem}{Theorem}
\newtheorem{lemma}[theorem]{Lemma}
\newtheorem{proposition}[theorem]{Proposition}
\newtheorem{corollary}[theorem]{Corollary}
\theoremstyle{definition}
\newtheorem{definition}[theorem]{Definition}
\theoremstyle{remark}
\newtheorem{remark}[theorem]{Remark}

\newcommand{\cube}{\{0,1\}}
\newcommand{\dist}{d_{\mathrm H}}
\newcommand{\supp}{\operatorname{supp}}
\newcommand{\divg}{\operatorname{div}}
\newcommand{\ov}{\operatorname{ov}}

\title{A Fixed-Backbone Exclusion for Length-70 Binary
  $(9,1)$ Covering Sequences}
\author{Ruturaj R Raval\\
  Independent Researcher\\
  \href{https://orcid.org/0000-0003-4930-8981}
  {ORCID: 0000-0003-4930-8981}}
\date{September 2026}

\begin{document}
\maketitle

\begin{abstract}
A binary $(n,R)$ covering sequence is a cyclic binary word whose
length-$n$ windows cover the Hamming cube within radius $R$.  The known
bounds for the minimum length at $(n,R)=(9,1)$ are
$62\leq L(9,1)\leq 71$.  We fix an explicit 64-edge balanced covering
support $B$ in the order-8 binary de Bruijn digraph.  First, a finite
graph-theoretic certificate proves that every connected nonnegative
integral circulation of total mass 70 uses at most 61 distinct edges of
$B$.  Second, an exhaustive path-cycle decomposition classifies the
exact-overlap-61 shell.  Among all $\binom{64}{3}=41{,}664$ triples of
omitted backbone edges, 188 residual flows survive and exactly eight
produce connected 70-edge circulations.  Six of these leave nine binary
9-words uncovered and two leave ten, so none is a radius-1 cover.
Consequently, every length-70 binary $(9,1)$ covering sequence, if one
exists, uses at most 60 distinct edges of $B$.  The overlap-61
classification is reproduced by separate Python and C++ implementations;
the prerequisite overlap theorem retains all 168 boundary residuals and
a semantic validator.  This result does not construct or exclude a
length-70 covering sequence and does not change the global bounds on
$L(9,1)$.
\end{abstract}

\section{Introduction}

Chung and Cooper introduced covering de Bruijn cycles as cyclic analogues
of covering codes \cite{chung-cooper}.  Chee, Etzion, Ta, and Vu later
developed general constructions for covering sequences and arrays
\cite{chee-et-al}.  For the binary radius-1 problem at window length 9,
their table gave
\[
  62\leq L(9,1)\leq 93.
\]
Rosin subsequently reported a 71-bit construction \cite{rosin}, giving
the current public bound
\[
  62\leq L(9,1)\leq 71.
\]

The present work does not improve either endpoint.  Instead, it gives an
exact finite restriction on any possible length-70 construction.  We fix
a 64-edge support $B$ that is balanced and covers all 512 binary 9-words
within radius 1, but has two directed-cycle components of sizes 4 and 60.
The support arose during exploratory search as the intersection of the
reported 71-edge cycle support and a disconnected 70-edge covering
support.  The theorem below depends only on the explicit fixed set $B$,
not on that discovery route.

There are two parts.  The first proves that a length-70 cyclic word can
use at most 61 distinct edges of $B$.  The cases of 64 and 63 retained
edges are excluded by short connector and detour bounds.  The case of 62
retained edges is closed by an exhaustive enumeration of 168 residual
flows.  The second part classifies the boundary value 61.  Exactly eight
connected circulations attain it, and none covers the Hamming cube.

\begin{table}[ht]
\centering
\begin{tabular}{@{}lll@{}}
\toprule
Source & Global bound & Contribution \\
\midrule
Chee et al.\ \cite{chee-et-al} & $62\leq L(9,1)\leq 93$ &
General constructions and table \\
Rosin \cite{rosin} & $62\leq L(9,1)\leq 71$ &
Explicit 71-bit construction \\
This work & unchanged &
Every length-70 cover has $\ov_B\leq 60$ \\
\bottomrule
\end{tabular}
\caption{Global bounds and the scoped result proved here.}
\label{tab:bounds}
\end{table}

\section{Covering sequences and de Bruijn circulations}

\begin{definition}
Let $x=x_0x_1\ldots x_{L-1}$ be a cyclic binary word.  Its length-$n$
window at position $i$ is
\[
  W_i=x_ix_{i+1}\cdots x_{i+n-1},
\]
with indices reduced modulo $L$.  The word $x$ is an $(n,R)$ covering
sequence if every $y\in\cube^n$ satisfies
\[
  \min_{0\leq i<L}\dist(y,W_i)\leq R.
\]
The minimum possible $L$ is denoted by $L(n,R)$.
\end{definition}

Let $D_8$ be the binary de Bruijn digraph with vertex set $\cube^8$ and
edge set $\cube^9$.  A 9-bit word $w=w_1\cdots w_9$ is the edge from
$w_1\cdots w_8$ to $w_2\cdots w_9$.  We identify binary words with their
integer values from 0 to 511.

For an edge multiplicity vector
$m\in\mathbb Z_{\geq 0}^{512}$, write
\[
  |m|=\sum_e m_e,\qquad
  \supp(m)=\{e:m_e>0\}.
\]
Its divergence is outgoing multiplicity minus incoming multiplicity at
each vertex.  It is a circulation when $\divg(m)=0$.

\begin{lemma}[Cyclic-word equivalence]
\label{lem:euler}
A nonzero vector $m\in\mathbb Z_{\geq0}^{512}$ is the window
multiplicity vector of a cyclic binary word if and only if $m$ is a
circulation and the underlying undirected graph of $\supp(m)$ is
connected.
\end{lemma}

\begin{proof}
A cyclic word orders its windows as one closed directed walk, so balance
and connectedness are immediate.  Conversely, replace every edge $e$ by
$m_e$ parallel copies.  A finite directed multigraph in which every
nonisolated vertex is balanced and the underlying graph is connected has
an Euler circuit.  Reading the appended bit along that circuit gives a
cyclic binary word.  Consecutive edges overlap in eight bits, and closure
of the circuit handles the wraparound windows.  Thus the resulting
window multiplicities are exactly $m$, including repeated edges and
periodic words.
\end{proof}

\section{The fixed backbone}

The backbone $B$ consists of the following 64 edges, written as decimal
values of 9-bit words:

\begin{center}
\small
\begin{tabular}{rrrrrrrr}
6&13&16&27&32&43&54&61\\
64&75&86&93&102&109&112&123\\
129&138&151&156&167&172&177&186\\
199&204&209&218&225&234&247&252\\
259&264&277&286&293&302&307&312\\
325&334&339&344&355&360&373&382\\
388&399&402&409&418&425&436&447\\
450&457&468&479&484&495&498&505
\end{tabular}
\end{center}

For a circulation $m$, define its distinct backbone overlap by
\[
  \ov_B(m)=|\supp(m)\cap B|.
\]

\begin{proposition}[Backbone properties]
\label{prop:backbone}
The support $B$ is balanced, has no loops, and covers all 512 binary
9-words within Hamming radius 1.  Its underlying graph has exactly two
components, each a directed cycle, with 4 and 60 edges respectively.
\end{proposition}

These properties are checked directly from the displayed edge list.  In
particular, they do not depend on a solver result.

\section{Residual-flow decomposition}

The finite classifications use the following standard decomposition.

\begin{lemma}[Integral path-cycle decomposition]
\label{lem:decomposition}
Let $a$ be a nonnegative integral edge flow in a finite directed
multigraph.  Suppose its positive divergence, counted with
multiplicity, has total value $d$.  Then $a$ is a sum of $d$ directed
source-to-sink walks and a nonnegative integral circulation.  The
circulation is a sum of directed closed walks.  Every walk in the
decomposition has length at most $|a|$.
\end{lemma}

\begin{proof}
Regard each unit of edge multiplicity as a separate edge copy.  Starting
at a positive-divergence vertex, follow unused outgoing copies until a
negative-divergence vertex is reached, deleting one unit along the
resulting directed walk.  At an intermediate vertex, failure to continue
would contradict its remaining divergence.  Repeating this operation
discharges all $d$ positive units.  The remaining multigraph is balanced.
Each nonempty weak component of a balanced directed multigraph has an
Euler circuit, giving the closed-walk decomposition.  No extracted walk
uses more edge copies than the original flow.
\end{proof}

\begin{lemma}[Completeness of bounded bit-walk enumeration]
\label{lem:enumeration}
Fix a residual mass $M$.  Enumerating every appended-bit string of
length $1,\ldots,M$ from each required source produces every directed
walk that can occur in Lemma~\ref{lem:decomposition}.  Enumerating every
cyclic binary word of length $1,\ldots,M$ and closing under sums produces
every balanced residual flow of mass at most $M$.  Pairing repeated
source and sink terminals in every distinct order therefore produces
every nonnegative integral residual flow of mass $M$ with the prescribed
divergence.
\end{lemma}

\begin{proof}
Every edge leaving a de Bruijn vertex is uniquely determined by its
appended bit, so a directed walk of length $\ell$ is represented by one
bit string of length $\ell$.  A directed closed walk of length $\ell$ in
$D_8$ has a cyclic appended-bit word of length $\ell$; this remains true
when $\ell<8$, in which case the associated 9-bit windows use the
periodic extension of the cyclic word.  Hence the closed-walk generator
contains every term in the balanced remainder of
Lemma~\ref{lem:decomposition}.  The recursive sum over all remaining
masses includes every multiset of closed walks.  Distinct pairings of
the source and sink multisets include repeated terminals, and set
deduplication removes only duplicate representations of the same flow.
\end{proof}

\section{Maximum possible backbone overlap}

The certificate establishes three finite facts.

\begin{proposition}[Finite connector facts]
\label{prop:finite}
For the backbone $B$:
\begin{enumerate}
\item A directed closed walk that starts in the 4-edge component,
touches the 60-edge component, and returns to its start has length at
least 7, even when backbone edges may be repeated.
\item If one backbone edge $e=u\to v$ is omitted, every directed
$u$-to-$v$ walk that avoids $e$ and touches the other backbone component
has length at least 10.
\item For every pair of omitted backbone edges, every residual flow of
mass 8 with the required divergence and avoiding the omitted edges gives
a disconnected combined support.
\end{enumerate}
\end{proposition}

For part 1, all
\[
  4(2^1+\cdots+2^6)=504
\]
appended-bit walks through length 6 are checked.  The bound is tight:
there are two length-7 connectors,
\[
\begin{split}
  &(205,411,310,108,217,435,358),\\
  &(306,100,201,403,294,76,153).
\end{split}
\]
For part 2, all
\[
  64(2^1+\cdots+2^7)=16{,}256
\]
walks through the length-7 residual budget are checked.  Breadth-first
search gives exact per-edge minima from 10 through 16, although only the
exclusion through 7 is needed below.

For part 3, Lemmas~\ref{lem:decomposition} and
\ref{lem:enumeration} reduce the problem to a finite enumeration.  The
certificate checks all $\binom{64}{2}=2{,}016$ omission pairs and retains
all 168 admissible residual flows.  Their combined support component
counts are
\[
\begin{array}{c|rrrr}
\text{components}&2&3&4&5\\ \hline
\text{residual flows}&36&76&50&6.
\end{array}
\]
No support is connected.  As a positive control, the same procedure at
residual mass 7 finds two connected completions of total mass 69.

\begin{theorem}[Exact overlap ceiling]
\label{thm:ceiling}
Every connected nonnegative integral circulation $m$ in $D_8$ with
$|m|=70$ satisfies
\[
  \ov_B(m)\leq 61.
\]
The bound is attained by a cyclic 70-bit word.
\end{theorem}

\begin{proof}
Suppose first that all 64 backbone edges occur.  Subtract one copy of
each edge of $B$.  The residual is a balanced flow of mass 6.
Connectivity of the combined support forces one residual component to
touch both components of $B$.  Its Euler circuit is a closed connector
of length at most 6, contradicting Proposition~\ref{prop:finite}(1).

Now suppose exactly one backbone edge $e=u\to v$ is omitted.  After
retaining one copy of each edge in $B\setminus\{e\}$, the residual has
mass 7 and divergence concentrated at $u$ and $v$.  Its component $H$
containing these terminals has a directed Euler trail from $u$ to $v$.
Indeed, the divergence sum of every weak residual component is zero, so
the unique positive and negative terminals must lie in the same
component.
Removing one edge from one directed-cycle component of $B$ leaves two
weak components.  Connectivity of the full support implies that some
residual component $K$ touches both of them.  This follows, for example,
from the connected bipartite incidence graph between backbone
components and residual components.

If $K=H$, the Euler trail is an $e$-avoiding detour through the other
backbone component of length at most 7, contradicting
Proposition~\ref{prop:finite}(2).  If $K\neq H$, then $K$ is balanced.
Since $H$ contains at least one edge, $K$ has mass at most 6, producing
a forbidden closed connector.

If exactly two backbone edges are omitted, the retained backbone part
has mass 62 and the residual has mass 8.  Proposition
\ref{prop:finite}(3) shows that every such completion is disconnected.
Thus overlaps 64, 63, and 62 are impossible.

Finally, the cyclic word
\[
\texttt{0010010111010100111000010000001100110000110110100010101100011110111111}
\]
has 70 distinct windows and overlap 61 with $B$.  It covers 503 of the
512 target words, so it certifies tightness of the graph-theoretic bound
but is not a covering sequence.
\end{proof}

\section{Classification of the overlap-61 shell}

Let $m$ have total mass 70 and $\ov_B(m)=61$.  It omits a triple
$E\subset B$.  Retain one copy of each edge in $B\setminus E$.  The
residual flow $a$ has mass 9, avoids $E$, and has divergence uniquely
determined by the retained part.

The source multiplicity is one, two, or three.  The enumeration
precomputes every directed walk of length at most 9 from each of the 64
backbone vertices:
\[
  64(2^1+\cdots+2^9)=65{,}408
\]
walks before endpoint grouping.  Balanced residual flows of masses
$0,\ldots,9$ occur in counts
\[
  1,2,4,8,16,32,64,128,256,512.
\]
Lemma~\ref{lem:enumeration} proves that the resulting classification is
exhaustive, including repeated windows and disconnected residual
remainders.

\begin{table}[ht]
\centering
\begin{tabular}{@{}r r r r@{}}
\toprule
Quantity & Count & Quantity & Count\\
\midrule
Omission triples & 41,664 & Active triples & 88\\
Raw decompositions & 192 & Distinct residual flows & 188\\
One component & 8 & Two components & 72\\
Three components & 80 & Four components & 28\\
Connected covers & 0 & Connected completions & 8\\
\bottomrule
\end{tabular}
\caption{Exact-overlap-61 enumeration.}
\label{tab:enumeration}
\end{table}

\begin{theorem}[Boundary-shell classification]
\label{thm:shell}
Exactly eight connected edge-multiplicity vectors of total mass 70 have
backbone overlap 61.  All eight have 70 distinct edges.  Six leave nine
binary 9-words uncovered within radius 1, and two leave ten.  In
particular, no length-70 binary $(9,1)$ covering sequence has overlap 61
with $B$.
\end{theorem}

\begin{proof}
For each of the $\binom{64}{3}=41{,}664$ omission triples, the algorithm
constructs every residual flow described by
Lemmas~\ref{lem:decomposition} and \ref{lem:enumeration}, rejects flows
using an omitted edge, and deduplicates the edge-multiplicity vectors.
Exactly 188 residuals remain.  Direct balance and connectivity checks
leave the eight vectors in Table~\ref{tab:completions}.  Direct Hamming
coverage checks give the listed uncovered sets, none empty.
\end{proof}

\begin{longtable}{@{}>{\raggedright\arraybackslash}p{0.13\textwidth}
  >{\raggedright\arraybackslash}p{0.38\textwidth}
  >{\raggedright\arraybackslash}p{0.38\textwidth}@{}}
\caption{The eight connected overlap-61 completions.  All residual edge
multiplicities are one.  Integers encode binary 9-words.}
\label{tab:completions}\\
\toprule
Omitted $E$ & Residual edges & Uncovered words\\
\midrule
\endfirsthead
\toprule
Omitted $E$ & Residual edges & Uncovered words\\
\midrule
\endhead
13, 307, 409 &
12, 25, 51, 97, 195, 269, 304, 390, 408 &
5, 15, 141, 311, 315, 401, 411, 413, 435\\
61, 102, 204 &
60, 103, 121, 207, 243, 317, 414, 460, 486 &
53, 63, 76, 98, 100, 110, 189, 196, 200\\
102, 151, 307 &
89, 101, 150, 179, 203, 300, 306, 358, 407 &
23, 98, 110, 149, 159, 230, 305, 311, 315\\
102, 204, 498 &
103, 121, 207, 242, 316, 414, 460, 486, 499 &
76, 98, 100, 110, 196, 200, 370, 496, 506\\
102, 307, 344 &
88, 101, 179, 203, 300, 306, 345, 358, 406 &
98, 110, 230, 305, 311, 315, 336, 346, 348, 472\\
167, 204, 409 &
105, 153, 166, 205, 211, 308, 332, 410, 423 &
39, 163, 165, 175, 196, 200, 206, 281, 401, 413\\
204, 360, 409 &
104, 153, 205, 211, 308, 332, 361, 410, 422 &
196, 200, 206, 281, 352, 362, 401, 413, 488\\
307, 409, 450 &
25, 51, 97, 194, 268, 304, 390, 408, 451 &
311, 315, 322, 401, 411, 413, 435, 448, 458\\
\bottomrule
\end{longtable}

\begin{corollary}
\label{cor:cover}
Every length-70 binary $(9,1)$ covering sequence, if one exists,
satisfies
\[
  \ov_B\leq 60.
\]
\end{corollary}

\begin{proof}
Theorem~\ref{thm:ceiling} gives overlap at most 61 for every cyclic word,
and Theorem~\ref{thm:shell} excludes overlap 61 for a cover.
\end{proof}

\section{Verification architecture}

The two stages have deliberately different verification claims.

\begin{table}[ht]
\centering
\begin{tabular}{@{}p{0.25\textwidth}p{0.65\textwidth}@{}}
\toprule
Stage & Checks\\
\midrule
Overlap ceiling &
The Python analyzer retains all 168 overlap-62 residual vectors.
A separate semantic validator checks every vector, reruns the 504
connector and 16,256 detour exclusions, reconstructs all histograms,
and checks the tight witness.\\
Overlap-61 shell &
Separate Python and C++ implementations agree on all aggregate counts
and all eight connected vectors.  A semantic validator checks all 188
residuals and every uncovered-word set.\\
Finite oracles &
Direct enumeration of every cyclic word in four smaller instances
checks zero-divergence, nonzero-divergence, three-omission, and
residual-mass-9 cases.\\
Integrity &
SHA-256 manifests authenticate the retained source, inputs, and
outputs.  Mutation tests reject missing or altered residuals.\\
\bottomrule
\end{tabular}
\caption{Verification layers.  Separate implementations improve
reproducibility but are not a substitute for external mathematical
review.}
\label{tab:verification}
\end{table}

The focused certificate sources use only the Python standard library;
the second overlap-61 implementation requires a C++20 compiler.  The
ancillary source archive contains deterministic manifests and exact
replay commands.

\section{Scope and future work}

The result is local to the explicit backbone $B$ and length 70.  It does
not show that $B$ is canonical, optimal among backbones, or useful at
other lengths.  It does not exclude a length-70 cover with overlap at
most 60.  Because covering-sequence existence is not generally monotone
in the length \cite{chung-cooper}, excluding length 70 alone would also
not establish a lower bound of 71 without additional arguments for the
shorter unresolved lengths.

The immediate use is exact-search reduction: any length-70 candidate may
be constrained to omit at least four edges of $B$.  More broadly, the
path-cycle method provides a certificate format for closing successive
overlap shells or comparing other disconnected covering supports.
Potential next steps are:
\begin{enumerate}
\item classify part or all of the overlap-60 shell;
\item search for a second backbone whose overlap restriction intersects
the present one strongly;
\item derive analytic inequalities that imply several shell exclusions
at once; and
\item continue independently verified construction search in the
remaining region.
\end{enumerate}

\section*{Data and code availability}

The code, explicit backbone, all retained residual flows, both
overlap-61 implementations, semantic validators, tests, and deterministic
ancillary archive are available from the public repository
\url{https://github.com/ruturajr-raval/binary-covering-sequence-9-1}.
Archived releases are available under the concept DOI
\href{https://doi.org/10.5281/zenodo.22260691}
{10.5281/zenodo.22260691}.  The release and ancillary source archive
contain the versioned publication evidence directories.

\section*{Acknowledgements}

The author thanks Christopher D. Rosin for making the 71-bit construction
and its provenance available, and the authors of the cited covering
sequence literature for their public manuscripts.

\begin{thebibliography}{9}

\bibitem{chung-cooper}
F.~R. K. Chung and J. Cooper,
\emph{De Bruijn cycles for covering codes},
arXiv:math/0310385, 2003.

\bibitem{chee-et-al}
Y.~M. Chee, T. Etzion, H. Ta, and V.~K. Vu,
\emph{Constructions of covering sequences and arrays},
arXiv:2502.08424, 2025.

\bibitem{rosin}
C.~D. Rosin,
\emph{New covering sequences based on a two-stage approach using large
language models and evolutionary algorithms},
arXiv:2505.23881, 2025.

\end{thebibliography}

\end{document}
